%-*-latex-*-

\documentclass[a4paper,11pt]{article}

% Language and fonts
%
\usepackage[utf8]{inputenc}
\usepackage[T1]{fontenc}
\usepackage[english]{babel}
\usepackage[hyphens]{url}
\usepackage{hyphenat}
\usepackage{array}
\usepackage{microtype}
\usepackage{ebgaramond}

\title{Tutoring notes}
\date{\oldstylenums{2026}}
\author{Christian Rinderknecht}

\frenchspacing  % Follow French conventions after a period

\begin{document}

\maketitle

\begin{itemize}

  \item \emph{Hogy vagy?} `How are you?' inform. --- \emph{Hogy van?}
    formal.

  \item \emph{Megvagyok.} `I am all right.' (neutral)

  \item \emph{Rosszul/Beteg vagyok.} `I am ill/sick.' (\emph{rossz}
    `bad')

  \item \emph{ön}: pers. pron. `you' sg. form.; \emph{önök} `you'
    pl. form.

  \item \emph{Mi újság?} `What's up?' inform.

  \item \emph{És veled?} `And what with you?' (\emph{vel-} `with' +
    \emph{-ed} 2nd pers. sg. `you' $\rightarrow$
    \emph{veled}). \emph{velem} `with me', \emph{veled} `with you
    (sg. inform.)', \emph{vele} `with him/her/it', \emph{velünk} `with
    us', \emph{veletek} `with you (pl. inform.)', \emph{velük} `with
    them'.

  \item \emph{Semmi különös/extra.} `Nothing special.'; \emph{Minden
  rendben.} `Everything is fine.'

  \item \emph{Milyen idő van?} or \emph{Milyen az idő?} `What's the
    weather?'; \emph{milyen} `what kind of'. Compare to \emph{melyik}
    `which (one) of'.

  \item \emph{Jól/Rossz idő van.} `The weather is good/bad.'

  \item \emph{Szeles idő van} `The weather is windy.'  \emph{szél}
    `wind'. \emph{Fúj a szél.} `The wind is blowing.'

  \item \emph{Esik az eső.} `It's raining.'; \emph{Esik a hó.}  `It's
    snowing.; \emph{Esős idő van.} `The weather is rainy.'

  \item \emph{Hideg van.} `It is cold.' or \emph{Hideg idő van.}  `The
    weather is cold.'; \emph{Felhő\textbf{s} idő van.} `The weather is
    cloudy.'; \emph{Hűvös idő van.} `The weather is chilly.'
    \emph{Kicsit beteg vagyok, mert hűvös az idő Egerben.}
    (\emph{mert}: `because')

  \item \emph{Fázom.} `I am/feel cold.' (vb. ind. \emph{fázik} `to be
    cold'). \emph{Megfáztam.} `I caught a cold.' \emph{Meg vagyok
    fázva.} `I have a cold.'

  \item \emph{Süt a nap.} `The sun shines.'; \emph{Napos idő van.}
    `The weather is sunny.'; \emph{Tiszta az ég.} `Clear skies.'

  \item \emph{Hány fok van?} `How many degrees is it?'

\end{itemize}

\newpage

\noindent Example of small talk:

\begin{itemize}

  \item[---] \emph{Szia! Hogy vagy?}

  \item[---] \emph{Szia! Jól vagyok, csak egy kicsit fáradt vagyok. És
  te?}

  \item[---] \emph{Megvagyok, köszönöm. Mi újság?}

  \item[---] \emph{Semmi különös. És veled?}

  \item[---] \emph{Velem sincs semmi extra. Egerben milyen idő van?}

  \item[---] \emph{Itt hűvös az idő. Budapesten hány fok van?}

\end{itemize}

\noindent Resuming the lesson:
\begin{itemize}

  \item \emph{Kérdezhetek [...]?} `May/Can I ask [...]?'
    (\emph{kérdez} vb. `to ask a question' + \emph{-het} potential +
    \emph{-ek} 1st pers. sg. pres. ind.)

  \item \emph{valami} indef. pron. `something' (\emph{vala-}
    `some-/-ever' (undefined existence) + interr. pron. \emph{mi}
    `what')

  \item \emph{Mit jelent az, hogy [...]?}  `What does [...] mean?'
    (\emph{jelenteni} vb. inf. `signify', `mean', `report'
    $\rightarrow$ \emph{jelent} 3rd
    pers. sg. ind. pres. indef.). \emph{Ez jelent valamit.} `This
    means something.' \emph{Ez sok\textbf{at} jelent.} `This means a
    lot.'

  \item \emph{mondják} `they say/tell' or pass. `is said/told'
    (vb. \emph{mond} `to say/tell' + \emph{-ják} 3rd
    pers. pl. ind. pres. def.) \emph{Hogy mondják [...]?} `How to say
    [...]?'

  \item \emph{Világos minden?} `Is everything/all clear/understood?';
    \emph{világ} `light (arch.)' (also `world') + \emph{-os} denom.
    adj.-form. `-ish' (ornative: `endowed with') $\rightarrow$
    `bright', `understood'.

  \item \emph{fogalmam} `my idea'; \emph{fogalom} `concept, idea'+
    \emph{-am} 1st pers sg. single\hyp{}poss. `my'.

  \item \emph{sincs} `not ... either', `doesn't have either', `there
    isn't either' (\emph{sem} + \emph{nincs} $\rightarrow$
    \emph{sincs}). Use 3rd pers.

    \item \emph{Elnézést!} `Excuse-me!', `Sorry!'. \emph{el-}
      vb. pref. `away, past, over' + \emph{néz} vb. `to look'
      $\rightarrow$ `to overlook or forgive a fault'.
\end{itemize}

Some useful expressions:
\begin{itemize}

  \item \emph{Meg vagyok fázva.} `I have the flu.'; \emph{Megfáztam.}
    `I got a cold.'

  \item \emph{A számlát kérem.} `The bill, please.' (definite);
    \emph{A számlát kérem \textbf{szépen}.} sounds softer;
    \emph{Szeretnék egy kávét.} (polite).

  \item \emph{Egy kávét kérek.} `I would like a coffee.' (indefinite);

  \item \emph{Nekem is kell egy tea.} `I also need a tea.' (no
    accusative); \emph{nekem} `to/for me', \emph{neked} `to/for you
    (sg. inform.)', \emph{önnek} `to/for you (sg. form.)', \emph{neki}
    `to/for him/her/it', \emph{nekünk} `to/for us', \emph{nektek}
    `to/for you (pl. inform.)', \emph{önöknek} `to/for you
    (pl. formal)', \emph{nekik} `to/for them'.

  \item \emph{Nem annyi\textbf{ra} hideg.} `It's not that cold.'

  \item \emph{Menn\textbf{em} kell.} `I have to go.' (\emph{menni}
    vb. inf. `to go': \emph{megyek}, \emph{mész}, \emph{megy},
    \emph{megyünk}, \emph{mentek}, \emph{mennek})

  \item \emph{Tele vagyok.} `I am full.'; \emph{Elegem van.} `I have
    had enough, I am fed up.'

  \item \emph{Semmi baj.} `No harm done.' (After an apology.)

\end{itemize}

\noindent Countries, nationalities and languages.
\begin{enumerate}

  \item \textbf{Countries ending in \emph{-ország}}: nationality and
    language without ending, for example \emph{Franciaország},
    \emph{francia}, \emph{francia}. When an adverb, the language takes
    the suffixes \emph{-ul}/\emph{-ül}, depending on vowel harmony,
    e.g. \emph{Kicsit beszélek franci\textbf{ául}.} `I speak a bit of
    French.'

  \item \textbf{Countries ending in \emph{-ia}}: nationality and
    language without ending, e.g. \emph{Dánia} (Denmark), \emph{dán}
    (Danish citizen), \emph{dán} (Danish language). Exceptions:
    \emph{Bulgária}, \emph{bolgár}, \emph{bolgár}; \emph{Ausztria},
    \emph{osztrák}.

  \item \textbf{Nationalities and languages ending in \emph{-i}}:
    \emph{Amerika} (country) $\rightarrow$ \emph{amerika\emph{i}}
    (nationality); \emph{Kína} (country) $\rightarrow$
    \emph{kína\emph{i}} (nationality, language).

  \item \textbf{The rest}: \emph{Belgium}, \emph{belga}; \emph{Japán},
    \emph{japán}, \emph{japán}; \emph{Anglia}, \emph{angol},
    \emph{angol}.

\end{enumerate}

\noindent Use \emph{van} when referring to locations, or time or
emotions that are adverbs, e.g. \emph{Ő kórházban \textbf{van}?} `Is
(s)he at the hospital?'; \emph{Ő jól \textbf{van}?} `Is (s)he well?'

\begin{itemize}

  \item \emph{Egy pillanat.} `(just) A moment.'

  \item \emph{házi feladat} `homework'

  \item \emph{Van kérdés?} `Any questions?'

  \item \emph{Mennyibe kerül a kávé?} `How much does the coffee cost?'
    \emph{Hány forint a kávé?} `How many forints does the coffee
    cost?'

    \item The times of the day are:
    \begin{center}
      \begin{tabular}{>{\itshape}ll}
        hajnal   & `dawn' \\
        reggel   & `morning' \\
        délelőtt & `forenoon' \\
        dél      & `noon' \\
        délután  & `afternoon' \\
        este     & `evening' \\
        éjszaka  & `night' \\
        éjfél    & `midnight'
      \end{tabular}
    \end{center}
    Use those qualifiers before the time of the day, e.g. \emph{Reggel
    tíz óra.}

  \item \emph{kor} `age' (written language, official forms), or
    historical periods in speech, \emph{-kor} marks a time `at ...'

  \item Some professions: \emph{mérnök} `engineer', \emph{színész(nő)}
    `actor', \emph{énekes(nő)} `singer', \emph{politikus}
    `politician', \emph{szakács} `cook', \emph{séf} `chef',
    \emph{író(nő)} `writer', \emph{sofőr} `driver', \emph{pék}
    `baker', \emph{cukrász} `pastry chef', \emph{fogorvos} `dentist',
    \emph{fodrász} `hairdresser', \emph{pincér} `waiter',
    \emph{szerelő} `mechanic', \emph{asszistens} `assistant',
    \emph{tanár(nő)} `teacher', \emph{kasszás, pénztáros} `cashier',
    \emph{eladó} `salesman', \emph{irodai munkás} `office worker',
    \emph{ügyvéd} `lawyer', \emph{rendőr} `policeman', \emph{nővér}
    `nurse, older/religious sister', \emph{tűzoltó} `fireman',
    \emph{mentős} `paramedic', \emph{tenisz\-ező} `tennis player',
    \emph{informatikus} `computer scientist', \emph{háziasszony}
    `housewife', \emph{pedagógus} `educator'.

  \item \emph{nyugdíjas} `pensioner, retiree' (\emph{nyugalom} `rest,
    peace' or \emph{nyugszik} `to rest' $\rightarrow$ \emph{nyug-} +
    \emph{díj} `price, fee, reward' + \emph{-as} someone who
    possesses).

  \item \emph{Mit csinálsz?} `What do you (sg. inform.) do?'

  \item \emph{olyan} + adj., \emph{mint} + noun: `as' + adj. `as' +
    noun. E.g. \emph{Ági olyan nagy, mint Péter.} `Ági is as tall as
    Péter.'

  \item \emph{körülbelül} `approximately, about'; \emph{körül}
    post. `around'; \emph{A ház körül.} `Around the house.';
    \emph{Nyolc óra körül.} `Around eight o'clock.'; \emph{kör}
    `circle' is a back formation from \emph{körül}. \emph{belül}
    `inside' (emphasis on the border); \emph{A város\textbf{on}
    belül.} `Within the city.'

  \item \emph{\ldots, ugye?} `\ldots, right?'

  \item \emph{miért/mert} `why/because'.

  \item \emph{tényleg} `really'.

  \item \emph{beszél\textbf{get}} `to chat, converse'. The suffix
    \emph{-get} is traditionally described as frequentative or
    iterative, although the resulting meaning is not necessarily
    repetitive. Another example: \emph{hall} vb. `to hear' +
    \emph{-gat} $\rightarrow$ \emph{hallgat} `to listen, to remain
    quiet'.

  \item \emph{igaz/hamis} `true/false', `honest/fake'. \emph{Az
  állítás igaz vagy hamis.} `The statement is true or false.'

  \item \emph{néha} `sometimes'.

  \item \emph{A barátom} `my friend'; \emph{Az \textbf{én} barátom.}
    emph. `\emph{my} friend' (e.g. as opposed to someone else's
    friend)

  \item \emph{férj/feleség} `husband/wife'; pref. \emph{fél-} `half' +
    3rd pers. sg. single\hyp{}poss. \emph{-e} $\rightarrow$
    \emph{f\textbf{e}le} + \emph{-ség} `state' $\rightarrow$
    \emph{feleség}.

  \item \emph{pont} `point'; \emph{pont\textbf{os}} adj. `accurate';
    \emph{pontos\textbf{an}} adv. `precisely' (\emph{-an}:
    deadj. adv.-form. suffix).

  \item \emph{hölgy} `maiden'; \emph{asszony} `matron'.

\end{itemize}

\noindent Family members:
\begin{itemize}

  \item \emph{nagymama} `grandmother' $\rightarrow$ \emph{a
  nagymam\textbf{ám}} `my grandmother'; \\ \emph{nagypapa}
    `grandfather' $\rightarrow$ \emph{a nagypap\textbf{ám}} `my
    grandfather'.

  \item \emph{édesanya} (polite) `mother' (\emph{édes} `sweet' +
    \emph{anya} `mother') $\rightarrow$ \emph{a édesany\textbf{ám}}
    `my mother'; \emph{anyuka} (affectionate) `mother' $\rightarrow$
    \emph{a anyuk\textbf{ám}} `my mother'.

  \item \emph{édesapa} (polite) `father' $\rightarrow$ \emph{a
  édesap\textbf{ám}} `my father'; \emph{apuka} (affectionate) `father'
    $\rightarrow$ \emph{apuk\textbf{ám}} `my father'.

  \item \emph{lány} `girl, daughter' $\rightarrow$ \emph{a
  lány\textbf{om}} `my daughter'; \emph{fiú} `boy, son' $\rightarrow$
    \emph{a fi\textbf{am}} `my son'.

  \item \emph{testvér} `sibling' (\emph{test} `body' + \emph{vér}
    `blood'). \emph{Egy testvérem van.} `I have one sibling.'

  \item \emph{nővér} `elder sister' $\rightarrow$ \emph{a
  nővér\textbf{em}} `my elder sister'; \\ \emph{báty(a)} `elder
    brother' $\rightarrow$ \emph{a báty\textbf{ám}} `my elder
    brother'.

  \item \emph{húg} `younger sister' $\rightarrow$ \emph{a
  húg\textbf{om}} `my younger sister'; \emph{öcs} `younger brother'
    $\rightarrow$ \emph{a öcs\textbf{ém}} `my younger brother'.

\end{itemize}

\noindent Descriptives:
\begin{itemize}

  \item \emph{idős} adj. (polite) `elderly person or old animal';
    \emph{régi} adj. `old (object)'; \emph{öreg} `old
    (person)'. Anton. \emph{fiatal} adj. `young (person)'.

  \item \emph{szép} adj. `beautiful' (woman); \emph{jóképű} `handsome'
    (man) (\emph{kép} `image, face'). Anton. \emph{csúnya}
    adj. `ugly'.

  \item \emph{alacsony} adj. `short
    (person)'. Anton. \emph{magas} adj. `tall'.

  \item \emph{sovány} adj. `thin' (person). Anton. \emph{köver}
    adj. `fat', \emph{molett} `plus\hyp{}size'.

  \item \emph{kedves} adj, `kind'; \emph{barátságos}
    adj. `friendly'. \\ Anton. \emph{barátságtlan} `unfriendly'
    (\emph{barát} + \emph{-ság} + \emph{-tlan}).

  \item \emph{rokonszenves} adj. `likeable, sympathetic' (\emph{rokon}
    `relative, kin', obs. `similar' + reform. \emph{szenv} `feeling' +
    \emph{-es} denom. adj.-form.) Anton. \emph{ellenszenves} adj.
    `unpleasant'.

  \item \emph{vidám} adj. `cheerful'. Anton. \emph{szomorú} adj. `sad'

  \item \emph{nyugodt} adj. `calm'. Anton. \emph{ideges} `nervous'.

  \item \emph{tülrelmes} adj. `patient' (\emph{türelem} `patience' +
    \emph{-es} denom. adj.-form.). Anton. \emph{türelmetlen} adj.
    `impatient' (\emph{türelem} + \emph{-etlen} priv.).

  \item \emph{szorgalmas} adj. `diligent'. Anton. \emph{lusta}
    adj. `lazy, idle'.

  \item \emph{friss} adj. `fresh'. Anton. \emph{fáradt} adj. `tired'.

  \item \emph{jókedvű} adj. `cheerful, good\hyp{}natured' (\emph{jó}
    adj. `good' + \emph{kedv} `mood, disposition' + \emph{-ű}
    denom. adj.-form.). Anton. \emph{rosszkedvű}
    adj. `ill\hyp{}natured'.

  \item \emph{szemüveges} `wearing glasses'; \emph{bajszos} `with a
    mustache'; \emph{szakállas} adj. `bearded'.

  \item \emph{kopasz} adj.  `bald'; \emph{szőke} adj. `blonde';
    \emph{barna hajú} adj. `brown\hyp{}haired'; \emph{fekete hajú}
    `black\hyp{}haired'; \emph{ősz hajú} `grey\hyp{}haired';
    \emph{vörös hajú} `ginger, red\hyp{}haired'.

   \item \emph{kék szemű} `blue\hyp{}eyed'; \emph{zöld szemű}
     `green\hyp{}eyed'; \emph{barna szemű} `brown\hyp{}eyed'.

\end{itemize}

\end{document}
